%!TEX root =  main.tex


We first derive an upper bound on $T_{\SD}$ that does not involve the expectation operator. 
Under a mild assumption on the form of $f$, we have the following result.
\begin{restatable}{lem}{sd}
    \label{prop:sd}
    Let $f(t) = 1 - e^{-\alpha t}$ for some constant $\alpha \in \mathbb{R}^+$. We have that
    \begin{equation*}
        T_{\SD} > \frac{\frac{1}{\alpha}\ln(\alpha V+1) + V}{1-\frac{1}{2(\alpha V+1)}} \times R\ .
    \end{equation*}
\end{restatable}
% \sd*
\begin{proof}
    First, the form of $f(x)$ is monotone increasing in the domain of $\alpha$. Hence, the function $g(k) = \frac{k+V}{f(k)}$ has a global minimum, taken at its turning point (critical point theorem). In the following, we derive this turning point $g'(k) = 0$. Taking the derivative with respect to $k$ and substituting the formula of $f(x)$, we have
    \begin{equation*}
        g'(k) = 0\iff (k+V)f'(k)=f(k) \iff \alpha (k+V)e^{-\alpha k} = 1-e^{-\alpha k}\ .
    \end{equation*}
    Rearrange the last equation:
    \begin{equation*}
        (\alpha(k+V)+1)e^{-\alpha k} = 1\ .
    \end{equation*}
    We use the Lambert $W$ function to solve this equation. Specifically, the Lambert $W$ function is defined as the inverse function of $f(w) = we^{w}$, i.e.,
    \begin{equation*}
        z = we^{w} \iff w = W(z)\ .
    \end{equation*}
    With this, let $u \coloneqq \alpha(k+V) + 1$. Then, the equation above can be rewritten as $-ue^{-u}=-e^{-(\alpha V + 1)}$. So,
    \begin{equation*}
        u = -W(-e^{-(\alpha V+1)})\ ,
    \end{equation*}
    and since $k = \frac{u}{\alpha} - ( V + \frac{1}{\alpha})$, we have
    \begin{equation*}
        k^*=-\frac{1}{\alpha}W(-e^{-(\alpha V + 1)})-(V+\frac{1}{\alpha})\ .
    \end{equation*}
    For $z \in (-e^{-1},0)$, there are two real branches: $W_0(z) \in (-1,0)$ and $W_{-1}(z) \in (-\infty,-1]$. Here $z=-e^{-(\alpha V+1)} \in (-e^{-1},0)$ and $u = \alpha(k+V)+1\geq1$. If we took $W_0$, we would have $u=-W_0(z) \in (0,1)$, a contradiction. Therefore, we take $W_{-1}$. Let $m \coloneqq \alpha V+1$. The expansion for the $-1$ branch as $m \to \infty$ is
    \begin{equation*}
        W_{-1}(z) = \ln(-z) - \ln(-\ln(-z)) + \frac{\ln(-\ln(-z))}{\ln(-z)} + \cdots\ .
    \end{equation*}
    Substitute $z=-e^{-m}$, we have
    \begin{equation*}
        W_{-1}(-e^{-m})=-m - \ln{m} - \frac{\ln{m}}{m} + \mathcal{O}\left(\frac{(\ln{m})^2}{m^2}\right)\ .
    \end{equation*}
    Plugging this equation to the equation of $k^*$, we have
    \begin{align*}
        k^*&=\frac{1}{\alpha}\left[ \ln{m} + \frac{\ln{m}}{m} + \mathcal{O}\left(\frac{(\ln{m})^2}{m^2} \right) \right] \\
         &= \frac{1}{\alpha}\ln(\alpha V + 1) + \frac{1}{\alpha} \frac{\ln(\alpha V + 1)}{\alpha V + 1} + \mathcal{O}\left( \frac{(\ln(\alpha V)^2}{\alpha (\alpha V)^2} \right)\ .
    \end{align*}
    With this, we now bound the range of $k^*$. First, recall that $(\alpha(k+V)+1)e^{-\alpha k} = 1\ .$. We rearrange the equation of $k$ into
    \begin{equation*}
        \alpha k = \ln(\alpha(k+V) + 1)\ .
    \end{equation*}
    Let $y \coloneqq \alpha k$ and define $g(y) \coloneqq \ln(\alpha V + 1 +  y) - y$. Then, $g'(y) = \frac{1}{\alpha V + 1 + y} - 1 < 0$, so $g$ is strictly decreasing and has a unique zero $y^*$. Define $y_0 \coloneqq \ln(\alpha V + 1)$ and $y_1 \coloneqq \ln(2(\alpha V + 1))$. We now show that $y_0 < y^* < y_1$.

    At $y_0$, we have $g(y_0) = \ln(1+\frac{\ln(\alpha V + 1)}{\alpha V + 1}) > 0$. This implies $y^* > y_0$.

    At $y_1$, $g(y_1) = \ln(1/2 + \frac{\ln(2(\alpha V+1))}{2(\alpha V + 1)}$. For $x \geq 1$, $\ln(2x) \leq x$. Hence, $1/2+\ln(2x)/(2x) < 1$. Therefore, $g(y_1) < 0$ for $\alpha V> 0$. This implies $y^* < y_1$. Therefore, substitute $k = y / \alpha$, we have
    \begin{equation*}
        \frac{1}{\alpha}\ln(\alpha V + 1) < k^* < \frac{1}{\alpha}\ln(2(\alpha V+ 1))\ .
    \end{equation*}
    As such, we have a bound for $T_{\SD}$ as
    \begin{align*}
        T_{\SD} &= \frac{k^*+V}{f(k^*)}    > \frac{\frac{1}{\alpha}\ln(\alpha V + 1) + V}{f(\frac{1}{\alpha} \ln(2(\alpha V + 1)))} =  \frac{\frac{1}{\alpha}\ln(\alpha V+1) + V}{1-\frac{1}{2(\alpha V+1)}}\ . \hfill \qedhere 
    \end{align*}
\end{proof}


On the other hand, under PSD, the expected time for verification is
\begin{equation*}
\begin{aligned}
\textstyle T_{\PSD} &\coloneqq \min_t \mathbb{E}[\sum_{s=1}^N\max(V, t)] \\ &= \min_t \mathbb{E}[N]\max{(V, t)}{f(t)}\ .
\end{aligned}
\end{equation*}
While it is obvious from the two definitions that $T_{\PSD} < T_{\SD}$, we prove a stronger result:
\begin{restatable}{lem}{psd}
    \label{prop:psd}
    Similarly, let $f(t) = 1 - e^{-\alpha t}$ for some constant $\alpha \in \mathbb{R}^+$. We have that
    \begin{equation*}
        T_{\PSD} = \frac{RV}{1-e^{-\alpha V}}
    \end{equation*}
\end{restatable}


\begin{proof}
    since $f(k)$ is increasing, it is easy to observe that $k^* \geq V$. Consider that for $k \geq V$, we have $(\max(V, k))' \geq f'(k)$. Hence, $k^* = V$. Substitute the value, we have
    \begin{equation*}
        T_{\PSD} = \frac{\max(V, t)}{f(V)} = \frac{V}{1-e^{-\alpha V}}\ . \qedhere 
    \end{equation*}
\end{proof}


\ratio*
\begin{proof}
    Using \cref{prop:sd} and \cref{prop:psd}, 
    \begin{equation*}
        \begin{aligned}
        \frac{T_{\SD}}{T_{\PSD}} &= \frac{\frac{1}{\alpha}\ln(\alpha V + 1) + V}{V}\frac{1-e^{-\alpha V}}{1-\frac{1}{2(\alpha V + 1)}}\ .    
        \end{aligned}
    \end{equation*}
    We are interested in when $\frac{1-e^{-\alpha V}}{1-\frac{1}{2(\alpha V + 1)}} \geq 1$, i.e., $e^{-\alpha V} \leq \frac{1}{2(\alpha V+1)}$. Thus, we solve
    \begin{equation*}
        e^{-\alpha x} = \frac{1}{2(\alpha x+1)}\ , \alpha > 0, x > 0\ .
    \end{equation*}
    set $t = \alpha x + 1 \implies x = (t-1)/\alpha$. Then, $e^{-\alpha x} = e^{-(t-1)} = e^{1-t}$. So,
    \begin{equation*}
        2te^{1-t}=1 \iff te^{-t}=\frac{1}{2e}\ .
    \end{equation*}
    Rewrite $(-t)e^{-t}=-\frac{1}{2e}$ and use the definition of Lambert W function, $W(z)e^{W(z)}=z$:
    \begin{equation*}
        -t = W(-\frac{1}{2e}) \implies t = -W(\frac{1}{-2e})\ .
    \end{equation*}
    Therefore,
    \begin{equation*}
        x = \frac{-W(-\frac{1}{2e}) - 1}{\alpha}\ .
    \end{equation*}
    Therefore, substitute back $V = x$, we have that when $V > \frac{-W(-\frac{1}{2e}) - 1}{\alpha}$, $\frac{1-e^{-\alpha V}}{1-\frac{1}{2(\alpha V + 1)}} \geq 1$. Further substitute the numerical solution of $W_{-1}(-\frac{1}{2e}) \approx -2.68$, we have that, when $\alpha V \geq 1.68$, $\frac{T_{\SD}}{T_{\PSD}} > \frac{\frac{1}{\alpha}\ln(\alpha V+1)+V}{V} = 1+\frac{\ln(\alpha V+1)}{\alpha V}$. Since the function $y = \frac{\ln(x+1)}{x}$ is monotone decreasing, the minimum value is taken at $\alpha V \approx 1.68$, giving $1+\frac{\ln(\alpha V + 1)}{\alpha V} \approx 1 + 0.59 = 1.59$.
\end{proof}